\documentclass[11pt,a4paper]{article}
\usepackage[utf8]{inputenc}
\usepackage[spanish]{babel}
\usepackage{amsmath}
\usepackage{amssymb}
\usepackage{graphicx}
\usepackage{geometry}
\usepackage{caption}
\usepackage{subcaption}
\usepackage{booktabs}
\usepackage{hyperref}
\usepackage{float}
\geometry{top=2cm, bottom=2cm, left=2.2cm, right=2.2cm}

\title{\textbf{Reporte Integral de Adquisición, Calidad de Señal y Análisis PCA}}
\author{Laboratorio de Sistemas Dinámicos (LSD) - Sistema Ñandú}
\date{2026-08-25}

\begin{document}
\maketitle

\section{Metadatos y Condiciones Experimentales}
\begin{itemize}
    \item \textbf{Fecha de Sesión:} 2026-08-25
    \item \textbf{Sujeto de Prueba:} No especificado
    \item \textbf{Tensión de Baterías:} 8.30 V y 8.20 V
    \item \textbf{Electrodo de Referencia (Tierra):} Frente
    \item \textbf{Total de Pruebas Registradas:} 5
\end{itemize}

\subsection{Configuración de Canales y Zonas Anatómicas}
\begin{table}[H]
\centering
\begin{tabular}{clp{9cm}}
\toprule
\textbf{Canal} & \textbf{Músculo Asignado} & \textbf{Detalle / Ubicación} \\
\midrule
Canal 0 & Canal 0 & Electrodo bipolar de superficie \\
Canal 1 & Canal 1 & Electrodo bipolar de superficie \\
Canal 2 & Canal 2 & Electrodo bipolar de superficie \\
Canal 3 & Canal 3 & Electrodo bipolar de superficie \\
\bottomrule
\end{tabular}
\caption{Distribución de canales y asignación muscular.}
\end{table}

\subsection{Registro Fotográfico del Montaje de Electrodos}
\begin{figure}[H]
\centering
\begin{subfigure}{0.45\textwidth}
    \includegraphics[width=\textwidth]{/home/santiago/repositorios/Nandu_SistemadeAdqusicionEMG/EMG_desarrollo/base_de_datos_electrodos/2026-08-25/O_Prueba1_Candela/photo.png}
    \caption{Montaje / Electrodo 1}
\end{subfigure}
\hfill
\begin{subfigure}{0.45\textwidth}
    \includegraphics[width=\textwidth]{/home/santiago/repositorios/Nandu_SistemadeAdqusicionEMG/EMG_desarrollo/base_de_datos_electrodos/2026-08-25/U_Prueba1_Candela/photo.png}
    \caption{Montaje / Electrodo 2}
\end{subfigure}
\vspace{0.3cm}
\begin{subfigure}{0.45\textwidth}
    \includegraphics[width=\textwidth]{/home/santiago/repositorios/Nandu_SistemadeAdqusicionEMG/EMG_desarrollo/base_de_datos_electrodos/2026-08-25/I_Prueba1_Candela/photo.png}
    \caption{Montaje / Electrodo 3}
\end{subfigure}
\hfill
\begin{subfigure}{0.45\textwidth}
    \includegraphics[width=\textwidth]{/home/santiago/repositorios/Nandu_SistemadeAdqusicionEMG/EMG_desarrollo/base_de_datos_electrodos/2026-08-25/E_Prueba1_Candela/photo.png}
    \caption{Montaje / Electrodo 4}
\end{subfigure}
\vspace{0.3cm}
\caption{Ubicación y fijación mecánica de los electrodos en la sesión.}
\end{figure}

\subsection{Armado y Fijación Mecánica}
Electrodos recortados con plancha de gel y doble cinta de refuerzo.

\subsection{Secuencia y Protocolo de Adquisición}
Protocolo de 10 repeticiones por vocal (A, E, I, O, U) y secuencias compuestas.

\subsection{Observaciones y Registro de Artefactos}
Picos involuntarios por deglución.

\newpage
\section{Análisis de Calidad de Señal y Relación Señal-Ruido}
Evaluación de los niveles de ruido interpulso, amplitudes pico y estabilidad temporal a lo largo de las pruebas registradas.

\newpage
\section{Análisis de Componentes Principales}
Proyecciones ortogonales y clustering para la discriminación de tareas vocálicas, evaluando el espacio completo de 3 canales y cada par muscular en 2D y 3D.

\subsection{Optimización de Hiperparámetros}
\begin{table}[H]
\centering
\begin{tabular}{lc}
\toprule
\textbf{Parámetro de Procesamiento} & \textbf{Valor Óptimo Hallado} \\
\midrule
Ventana de Suavizado ($\tau_{\text{suavizado}}$) & 125 ms \\
Puntos de Remuestreo ($W_{\text{ciclo}}$) & 20 muestras \\
Factor de Supresión de Ruido ($\alpha_{\text{ruido}}$) & 0.1 \\
Factor $Q$ del Filtro Notch & 2.0 \\
Exactitud Máxima de Clasificación & \textbf{83.66\%} \\
Coeficiente Silhouette & 0.1477 \\
\bottomrule
\end{tabular}
\caption{Hiperparámetros óptimos seleccionados mediante búsqueda en grilla.}
\end{table}

\newpage
\subsection{Tríada Completa (Canales 0, 1 y 2)}
\begin{figure}[H]
\centering
\begin{subfigure}{0.48\textwidth}
    \includegraphics[width=\textwidth]{/home/santiago/repositorios/Nandu_SistemadeAdqusicionEMG/reportes_experimentos/analisis_pca_generado/triada/heatmap_confusion_pca_2d.png}
    \caption{Matriz de Confusión 2D}
\end{subfigure}
\hfill
\begin{subfigure}{0.48\textwidth}
    \includegraphics[width=\textwidth]{/home/santiago/repositorios/Nandu_SistemadeAdqusicionEMG/reportes_experimentos/analisis_pca_generado/triada/heatmap_confusion_pca_3d.png}
    \caption{Matriz de Confusión 3D}
\end{subfigure}
\caption{Matrices de confusión y tasas de clasificación para Tríada Completa (Canales 0, 1 y 2).}
\end{figure}

\newpage
\subsection{Par de Canales 0 y 1}
\begin{figure}[H]
\centering
\begin{subfigure}{0.48\textwidth}
    \includegraphics[width=\textwidth]{/home/santiago/repositorios/Nandu_SistemadeAdqusicionEMG/reportes_experimentos/analisis_pca_generado/par_0_1/heatmap_confusion_pca_2d.png}
    \caption{Matriz de Confusión 2D}
\end{subfigure}
\hfill
\begin{subfigure}{0.48\textwidth}
    \includegraphics[width=\textwidth]{/home/santiago/repositorios/Nandu_SistemadeAdqusicionEMG/reportes_experimentos/analisis_pca_generado/par_0_1/heatmap_confusion_pca_3d.png}
    \caption{Matriz de Confusión 3D}
\end{subfigure}
\caption{Matrices de confusión y tasas de clasificación para Par de Canales 0 y 1.}
\end{figure}

\newpage
\subsection{Par de Canales 1 y 2}
\begin{figure}[H]
\centering
\begin{subfigure}{0.48\textwidth}
    \includegraphics[width=\textwidth]{/home/santiago/repositorios/Nandu_SistemadeAdqusicionEMG/reportes_experimentos/analisis_pca_generado/par_1_2/heatmap_confusion_pca_2d.png}
    \caption{Matriz de Confusión 2D}
\end{subfigure}
\hfill
\begin{subfigure}{0.48\textwidth}
    \includegraphics[width=\textwidth]{/home/santiago/repositorios/Nandu_SistemadeAdqusicionEMG/reportes_experimentos/analisis_pca_generado/par_1_2/heatmap_confusion_pca_3d.png}
    \caption{Matriz de Confusión 3D}
\end{subfigure}
\caption{Matrices de confusión y tasas de clasificación para Par de Canales 1 y 2.}
\end{figure}

\newpage
\subsection{Par de Canales 0 y 2}
\begin{figure}[H]
\centering
\begin{subfigure}{0.48\textwidth}
    \includegraphics[width=\textwidth]{/home/santiago/repositorios/Nandu_SistemadeAdqusicionEMG/reportes_experimentos/analisis_pca_generado/par_0_2/heatmap_confusion_pca_2d.png}
    \caption{Matriz de Confusión 2D}
\end{subfigure}
\hfill
\begin{subfigure}{0.48\textwidth}
    \includegraphics[width=\textwidth]{/home/santiago/repositorios/Nandu_SistemadeAdqusicionEMG/reportes_experimentos/analisis_pca_generado/par_0_2/heatmap_confusion_pca_3d.png}
    \caption{Matriz de Confusión 3D}
\end{subfigure}
\caption{Matrices de confusión y tasas de clasificación para Par de Canales 0 y 2.}
\end{figure}

\newpage
\section{Registros Detallados por Prueba}
A continuación se detallan individualmente todas las pruebas registradas en orden cronológico, indicando su hora exacta y tiempo relativo de sesión.

\newpage
\subsection{Prueba 1: O\_Prueba1\_Candela (14:19:55, $\Delta t = 0.0$ min)}
\begin{itemize}
    \item \textbf{Tarea / Fonema:} O
    \item \textbf{Hora de Registro:} 14:19:55
    \item \textbf{Tiempo Relativo de Sesión:} 0.0 minutos
\end{itemize}

\begin{figure}[H]
\centering
\includegraphics[width=0.65\textwidth]{/home/santiago/repositorios/Nandu_SistemadeAdqusicionEMG/EMG_desarrollo/base_de_datos_electrodos/2026-08-25/O_Prueba1_Candela/patron_muscular_grabacion.png}
\caption{Patrón de activación muscular promedio para la Prueba 1.}
\end{figure}

\begin{figure}[H]
\centering
\includegraphics[width=0.88\textwidth]{/home/santiago/repositorios/Nandu_SistemadeAdqusicionEMG/EMG_desarrollo/base_de_datos_electrodos/2026-08-25/O_Prueba1_Candela/plot_paper_combined.png}
\caption{Registro temporal y espectral combinado con ruido restado (Prueba 1).}
\end{figure}

\begin{figure}[H]
\centering
\includegraphics[width=0.88\textwidth]{/home/santiago/repositorios/Nandu_SistemadeAdqusicionEMG/EMG_desarrollo/base_de_datos_electrodos/2026-08-25/O_Prueba1_Candela/plot_calibrado_2026-08-25_o_Prueba1_Candela.png}
\caption{Señales calibradas en microvoltios (filtro notch 50 Hz, pasabanda 20--500 Hz y RMS 75 ms).}
\end{figure}

\begin{figure}[H]
\centering
\includegraphics[width=0.65\textwidth]{/home/santiago/repositorios/Nandu_SistemadeAdqusicionEMG/EMG_desarrollo/base_de_datos_electrodos/2026-08-25/O_Prueba1_Candela/canal_0/avg_lider.png}
\vspace{0.3cm}
\caption{Canal 0: Canal 0 (Prueba 1).}
\end{figure}

\begin{figure}[H]
\centering
\includegraphics[width=0.65\textwidth]{/home/santiago/repositorios/Nandu_SistemadeAdqusicionEMG/EMG_desarrollo/base_de_datos_electrodos/2026-08-25/O_Prueba1_Candela/canal_1/avg_lider.png}
\vspace{0.3cm}
\caption{Canal 1: Canal 1 (Prueba 1).}
\end{figure}

\begin{figure}[H]
\centering
\includegraphics[width=0.65\textwidth]{/home/santiago/repositorios/Nandu_SistemadeAdqusicionEMG/EMG_desarrollo/base_de_datos_electrodos/2026-08-25/O_Prueba1_Candela/canal_2/avg_lider.png}
\vspace{0.3cm}
\caption{Canal 2: Canal 2 (Prueba 1).}
\end{figure}

\newpage
\subsection{Prueba 2: U\_Prueba1\_Candela (14:20:39, $\Delta t = 0.7$ min)}
\begin{itemize}
    \item \textbf{Tarea / Fonema:} U
    \item \textbf{Hora de Registro:} 14:20:39
    \item \textbf{Tiempo Relativo de Sesión:} 0.7 minutos
\end{itemize}

\begin{figure}[H]
\centering
\includegraphics[width=0.65\textwidth]{/home/santiago/repositorios/Nandu_SistemadeAdqusicionEMG/EMG_desarrollo/base_de_datos_electrodos/2026-08-25/U_Prueba1_Candela/patron_muscular_grabacion.png}
\caption{Patrón de activación muscular promedio para la Prueba 2.}
\end{figure}

\begin{figure}[H]
\centering
\includegraphics[width=0.88\textwidth]{/home/santiago/repositorios/Nandu_SistemadeAdqusicionEMG/EMG_desarrollo/base_de_datos_electrodos/2026-08-25/U_Prueba1_Candela/plot_paper_combined.png}
\caption{Registro temporal y espectral combinado con ruido restado (Prueba 2).}
\end{figure}

\begin{figure}[H]
\centering
\includegraphics[width=0.88\textwidth]{/home/santiago/repositorios/Nandu_SistemadeAdqusicionEMG/EMG_desarrollo/base_de_datos_electrodos/2026-08-25/U_Prueba1_Candela/plot_calibrado_2026-08-25_u_Prueba1_Candela.png}
\caption{Señales calibradas en microvoltios (filtro notch 50 Hz, pasabanda 20--500 Hz y RMS 75 ms).}
\end{figure}

\begin{figure}[H]
\centering
\includegraphics[width=0.65\textwidth]{/home/santiago/repositorios/Nandu_SistemadeAdqusicionEMG/EMG_desarrollo/base_de_datos_electrodos/2026-08-25/U_Prueba1_Candela/canal_0/avg_lider.png}
\vspace{0.3cm}
\caption{Canal 0: Canal 0 (Prueba 2).}
\end{figure}

\begin{figure}[H]
\centering
\includegraphics[width=0.65\textwidth]{/home/santiago/repositorios/Nandu_SistemadeAdqusicionEMG/EMG_desarrollo/base_de_datos_electrodos/2026-08-25/U_Prueba1_Candela/canal_1/avg_lider.png}
\vspace{0.3cm}
\caption{Canal 1: Canal 1 (Prueba 2).}
\end{figure}

\begin{figure}[H]
\centering
\includegraphics[width=0.65\textwidth]{/home/santiago/repositorios/Nandu_SistemadeAdqusicionEMG/EMG_desarrollo/base_de_datos_electrodos/2026-08-25/U_Prueba1_Candela/canal_2/avg_lider.png}
\vspace{0.3cm}
\caption{Canal 2: Canal 2 (Prueba 2).}
\end{figure}

\newpage
\subsection{Prueba 3: I\_Prueba1\_Candela (14:23:36, $\Delta t = 3.7$ min)}
\begin{itemize}
    \item \textbf{Tarea / Fonema:} I
    \item \textbf{Hora de Registro:} 14:23:36
    \item \textbf{Tiempo Relativo de Sesión:} 3.7 minutos
\end{itemize}

\begin{figure}[H]
\centering
\includegraphics[width=0.65\textwidth]{/home/santiago/repositorios/Nandu_SistemadeAdqusicionEMG/EMG_desarrollo/base_de_datos_electrodos/2026-08-25/I_Prueba1_Candela/patron_muscular_grabacion.png}
\caption{Patrón de activación muscular promedio para la Prueba 3.}
\end{figure}

\begin{figure}[H]
\centering
\includegraphics[width=0.88\textwidth]{/home/santiago/repositorios/Nandu_SistemadeAdqusicionEMG/EMG_desarrollo/base_de_datos_electrodos/2026-08-25/I_Prueba1_Candela/plot_paper_combined.png}
\caption{Registro temporal y espectral combinado con ruido restado (Prueba 3).}
\end{figure}

\begin{figure}[H]
\centering
\includegraphics[width=0.88\textwidth]{/home/santiago/repositorios/Nandu_SistemadeAdqusicionEMG/EMG_desarrollo/base_de_datos_electrodos/2026-08-25/I_Prueba1_Candela/plot_calibrado_2026-08-25_i_Prueba1_Candela.png}
\caption{Señales calibradas en microvoltios (filtro notch 50 Hz, pasabanda 20--500 Hz y RMS 75 ms).}
\end{figure}

\begin{figure}[H]
\centering
\includegraphics[width=0.65\textwidth]{/home/santiago/repositorios/Nandu_SistemadeAdqusicionEMG/EMG_desarrollo/base_de_datos_electrodos/2026-08-25/I_Prueba1_Candela/canal_0/avg_lider.png}
\vspace{0.3cm}
\caption{Canal 0: Canal 0 (Prueba 3).}
\end{figure}

\begin{figure}[H]
\centering
\includegraphics[width=0.65\textwidth]{/home/santiago/repositorios/Nandu_SistemadeAdqusicionEMG/EMG_desarrollo/base_de_datos_electrodos/2026-08-25/I_Prueba1_Candela/canal_1/avg_lider.png}
\vspace{0.3cm}
\caption{Canal 1: Canal 1 (Prueba 3).}
\end{figure}

\begin{figure}[H]
\centering
\includegraphics[width=0.65\textwidth]{/home/santiago/repositorios/Nandu_SistemadeAdqusicionEMG/EMG_desarrollo/base_de_datos_electrodos/2026-08-25/I_Prueba1_Candela/canal_2/avg_lider.png}
\vspace{0.3cm}
\caption{Canal 2: Canal 2 (Prueba 3).}
\end{figure}

\newpage
\subsection{Prueba 4: E\_Prueba1\_Candela (14:24:20, $\Delta t = 4.4$ min)}
\begin{itemize}
    \item \textbf{Tarea / Fonema:} E
    \item \textbf{Hora de Registro:} 14:24:20
    \item \textbf{Tiempo Relativo de Sesión:} 4.4 minutos
\end{itemize}

\begin{figure}[H]
\centering
\includegraphics[width=0.65\textwidth]{/home/santiago/repositorios/Nandu_SistemadeAdqusicionEMG/EMG_desarrollo/base_de_datos_electrodos/2026-08-25/E_Prueba1_Candela/patron_muscular_grabacion.png}
\caption{Patrón de activación muscular promedio para la Prueba 4.}
\end{figure}

\begin{figure}[H]
\centering
\includegraphics[width=0.88\textwidth]{/home/santiago/repositorios/Nandu_SistemadeAdqusicionEMG/EMG_desarrollo/base_de_datos_electrodos/2026-08-25/E_Prueba1_Candela/plot_paper_combined.png}
\caption{Registro temporal y espectral combinado con ruido restado (Prueba 4).}
\end{figure}

\begin{figure}[H]
\centering
\includegraphics[width=0.88\textwidth]{/home/santiago/repositorios/Nandu_SistemadeAdqusicionEMG/EMG_desarrollo/base_de_datos_electrodos/2026-08-25/E_Prueba1_Candela/plot_calibrado_2026-08-25_e_Prueba1_Candela.png}
\caption{Señales calibradas en microvoltios (filtro notch 50 Hz, pasabanda 20--500 Hz y RMS 75 ms).}
\end{figure}

\begin{figure}[H]
\centering
\includegraphics[width=0.65\textwidth]{/home/santiago/repositorios/Nandu_SistemadeAdqusicionEMG/EMG_desarrollo/base_de_datos_electrodos/2026-08-25/E_Prueba1_Candela/canal_0/avg_lider.png}
\vspace{0.3cm}
\caption{Canal 0: Canal 0 (Prueba 4).}
\end{figure}

\begin{figure}[H]
\centering
\includegraphics[width=0.65\textwidth]{/home/santiago/repositorios/Nandu_SistemadeAdqusicionEMG/EMG_desarrollo/base_de_datos_electrodos/2026-08-25/E_Prueba1_Candela/canal_1/avg_lider.png}
\vspace{0.3cm}
\caption{Canal 1: Canal 1 (Prueba 4).}
\end{figure}

\begin{figure}[H]
\centering
\includegraphics[width=0.65\textwidth]{/home/santiago/repositorios/Nandu_SistemadeAdqusicionEMG/EMG_desarrollo/base_de_datos_electrodos/2026-08-25/E_Prueba1_Candela/canal_2/avg_lider.png}
\vspace{0.3cm}
\caption{Canal 2: Canal 2 (Prueba 4).}
\end{figure}

\newpage
\subsection{Prueba 5: A\_Prueba1\_Candela (14:25:04, $\Delta t = 5.2$ min)}
\begin{itemize}
    \item \textbf{Tarea / Fonema:} A
    \item \textbf{Hora de Registro:} 14:25:04
    \item \textbf{Tiempo Relativo de Sesión:} 5.2 minutos
\end{itemize}

\begin{figure}[H]
\centering
\includegraphics[width=0.65\textwidth]{/home/santiago/repositorios/Nandu_SistemadeAdqusicionEMG/EMG_desarrollo/base_de_datos_electrodos/2026-08-25/A_Prueba1_Candela/patron_muscular_grabacion.png}
\caption{Patrón de activación muscular promedio para la Prueba 5.}
\end{figure}

\begin{figure}[H]
\centering
\includegraphics[width=0.88\textwidth]{/home/santiago/repositorios/Nandu_SistemadeAdqusicionEMG/EMG_desarrollo/base_de_datos_electrodos/2026-08-25/A_Prueba1_Candela/plot_paper_combined.png}
\caption{Registro temporal y espectral combinado con ruido restado (Prueba 5).}
\end{figure}

\begin{figure}[H]
\centering
\includegraphics[width=0.88\textwidth]{/home/santiago/repositorios/Nandu_SistemadeAdqusicionEMG/EMG_desarrollo/base_de_datos_electrodos/2026-08-25/A_Prueba1_Candela/plot_calibrado_2026-08-25_a_Prueba1_Candela.png}
\caption{Señales calibradas en microvoltios (filtro notch 50 Hz, pasabanda 20--500 Hz y RMS 75 ms).}
\end{figure}

\begin{figure}[H]
\centering
\includegraphics[width=0.65\textwidth]{/home/santiago/repositorios/Nandu_SistemadeAdqusicionEMG/EMG_desarrollo/base_de_datos_electrodos/2026-08-25/A_Prueba1_Candela/canal_0/avg_lider.png}
\vspace{0.3cm}
\caption{Canal 0: Canal 0 (Prueba 5).}
\end{figure}

\begin{figure}[H]
\centering
\includegraphics[width=0.65\textwidth]{/home/santiago/repositorios/Nandu_SistemadeAdqusicionEMG/EMG_desarrollo/base_de_datos_electrodos/2026-08-25/A_Prueba1_Candela/canal_1/avg_lider.png}
\vspace{0.3cm}
\caption{Canal 1: Canal 1 (Prueba 5).}
\end{figure}

\begin{figure}[H]
\centering
\includegraphics[width=0.65\textwidth]{/home/santiago/repositorios/Nandu_SistemadeAdqusicionEMG/EMG_desarrollo/base_de_datos_electrodos/2026-08-25/A_Prueba1_Candela/canal_2/avg_lider.png}
\vspace{0.3cm}
\caption{Canal 2: Canal 2 (Prueba 5).}
\end{figure}

\end{document}